\emph{Two explicit completion checks for the preceding argument.}
The flag description can be made completely concrete here. The only
strict flags with three strata are
$p<S<U$, where $p$ is a boundary crossing and $S$ is a
thrice-punctured boundary component. The normal circle bundle over
$S$ is topologically trivial, since $H^2(S,\mathbb Z)=0$.
Consequently its link group is $G_S\times\mathbb Z$ and its map to
$G_S$ is the projection. The other adjacent link has group
$\mathbb Z$, mapping to $G_S$ by the relevant puncture loop. The
three-stratum flag space is their homotopy pullback, hence has group
\[
 \mathbb Z\mathbin{\times}_{G_S}(G_S\times\mathbb Z)
                         \simeq\mathbb Z^2.
\]
After profinite completion the same identity holds, with hats: the
projection is split and all these products carry the full induced
profinite topology. Thus the finitely many nontrivial Segal gluing
identities survive completion in this geometry. This supplies the
hypothesis needed to apply the spatial and profinite decollage
equivalences of \cite[Theorem~2.7.4 and equivalence~(2.8.9)]{BGH2020};
no preservation of arbitrary homotopy pullbacks by completion is
asserted.

It also verifies universality for \emph{strict} finite categorical
quotients, which does not follow from density alone. Choose one point
in each stratum. For each incident pair $S<T$, the link is connected,
so its Hom set is a transitive $(G_T,G_S)$-biset. Upon choosing one
arrow its stabilizer is the image of its link group $H_{S,T}$ in
$G_T\times G_S$ (using the inverse right action). The injectivity
already proved gives the completed Hom space
\[
  (\widehat G_T\times\widehat G_S)
           /\operatorname{im}(\widehat H_{S,T}).
\]
The denominator is compact and therefore closed. A functor from the
discrete exit category to a finite category restricts to homomorphisms
from each $G_S$ to finite automorphism groups; these extend to
$\widehat G_S$. Its map on an incident Hom set is determined by the
image of one arrow and the finite group actions. The stabilizer
condition extends to the closure, so this map extends uniquely and
continuously to the displayed profinite Hom space. Composable pairs
of discrete arrows are dense in the corresponding product of Hom
spaces. Hence the extended maps preserve composition, since the
two resulting continuous maps to a finite discrete set agree on that
dense subset. This proves that every finite categorical quotient
extends, and hence that the compact category constructed from the
completed flags has precisely the strict finite quotients of $C$.
The context-congruence argument then identifies it with their inverse
limit.

Finally, the finite bound used for coherent representations has an
exact model-theoretic justification. The full coefficient category on
finitely many $\pi$-finite spaces has finitely many object classes,
finite homotopy groups in every mapping space, and a common finite
truncation bound on those mapping spaces. A minimal quasicategory
model is degreewise finite by Blom--Moerdijk, Lemma~A.10, and has a
finite coskeletality bound by Lurie, Proposition~2.3.4.18 and
Remark~2.3.4.4; this is exactly the conclusion recorded in
Blom--Moerdijk, Remark~A.11
\cite{FA_BlomMoerdijk2023,FA_LurieHTT2009}.
One takes this equivalent minimal model before using finite coherent
data; no strict coskeletality assertion about an arbitrary homotopy
coherent nerve is necessary. Thus a continuous coherent diagram is
determined by finitely many clopen data on products of Hom spaces.
The bound is taken separately for each finite-object diagram and for
each finite simplex of its mapping space, rather than for the whole
class of all $\pi$-finite spaces.
